% =====================================================================
%  Báo cáo đồ án môn Lập trình .NET — Đề tài GearGo
%  Nhóm 10 — GVHD: Thầy Nguyễn Thanh Sang
%  Trường Đại học Sài Gòn — Khoa Công nghệ Thông tin
%  Tháng 12 / 2026
%
%  Biên dịch: latexmk -xelatex main.tex   (khuyến nghị)
%             hoặc: xelatex main.tex ba lần
% =====================================================================
\documentclass[a4paper,13pt,oneside]{report}

% --- Khổ giấy & lề chuẩn báo cáo Việt Nam (trên 2 — dưới 2 — trái 3 — phải 2 cm)
\usepackage[a4paper,top=2cm,bottom=2cm,left=3cm,right=2cm]{geometry}

% --- Font tiếng Việt (XeLaTeX)
\usepackage{fontspec}
\usepackage{polyglossia}
\setdefaultlanguage{vietnamese}
\setmainfont{Times New Roman}
\setmonofont{DejaVu Sans Mono}[Scale=0.85]

% --- Khoảng cách dòng 1.5
\usepackage{setspace}
\onehalfspacing

% --- Các gói thường dùng
\usepackage{graphicx}
\usepackage{xcolor}
\usepackage{array}
\usepackage{tabularx}
\usepackage{longtable}
\usepackage{booktabs}
\usepackage{multirow}
\usepackage{enumitem}
\usepackage{float}
\usepackage{caption}
\usepackage{subcaption}
\usepackage{listings}
\usepackage{hyperref}
\usepackage{fancyhdr}
\usepackage{titlesec}
\usepackage{tocloft}

% --- Hyperref cấu hình
\hypersetup{
    colorlinks = true,
    linkcolor  = black,
    citecolor  = black,
    urlcolor   = blue,
    bookmarksnumbered = true,
    pdftitle   = {Báo cáo Đồ án GearGo — Nhóm 10},
    pdfauthor  = {Nhóm 10 — Khoa CNTT — Đại học Sài Gòn}
}

% --- Header & Footer
\pagestyle{fancy}
\fancyhf{}
\fancyhead[L]{\small GearGo --- Nền tảng cho thuê thiết bị dã ngoại}
\fancyhead[R]{\small Nhóm 10 --- Lập trình .NET}
\fancyfoot[C]{\thepage}
\renewcommand{\headrulewidth}{0.4pt}
\renewcommand{\footrulewidth}{0pt}

% --- Định dạng tiêu đề chương
\titleformat{\chapter}[hang]
  {\normalfont\Large\bfseries\MakeUppercase}
  {CHƯƠNG \thechapter.}{0.5em}{}
\titlespacing*{\chapter}{0pt}{-20pt}{20pt}

% --- Cấu hình code listing (C#, TypeScript, SQL)
\definecolor{codebg}{RGB}{245,245,245}
\definecolor{codekeyword}{RGB}{0,0,160}
\definecolor{codestring}{RGB}{160,0,0}
\definecolor{codecomment}{RGB}{0,120,0}

\lstdefinelanguage{csharp}{
  morekeywords={class,public,private,protected,internal,static,void,int,string,bool,
                async,await,Task,var,new,return,if,else,for,foreach,while,using,
                namespace,interface,readonly,override,virtual,abstract,sealed},
  sensitive=true,
  morecomment=[l]{//},
  morecomment=[s]{/*}{*/},
  morestring=[b]"
}

\lstset{
  backgroundcolor=\color{codebg},
  basicstyle=\ttfamily\footnotesize,
  keywordstyle=\color{codekeyword}\bfseries,
  stringstyle=\color{codestring},
  commentstyle=\color{codecomment}\itshape,
  breaklines=true,
  frame=single,
  framerule=0.3pt,
  showstringspaces=false,
  tabsize=2,
  captionpos=b,
  numbers=left,
  numberstyle=\tiny\color{gray},
  numbersep=8pt
}

% --- Định nghĩa cho bảng tabularx cột trung tâm
\newcolumntype{C}[1]{>{\centering\arraybackslash}p{#1}}

% =====================================================================
\begin{document}

% ---------- BÌA & MỤC LỤC ----------
\pagenumbering{roman}
% =====================================================================
%  BÌA BÁO CÁO — Đại học Sài Gòn — Khoa CNTT
% =====================================================================
\begin{titlepage}
\begin{center}

{\large\textbf{TRƯỜNG ĐẠI HỌC SÀI GÒN}}\\[0.3cm]
{\large\textbf{KHOA CÔNG NGHỆ THÔNG TIN}}\\[0.3cm]
{\large\textbf{HỌC PHẦN: CÔNG NGHỆ .NET}}\\[2cm]

% Logo trường (đặt file logo_sgu.png vào thư mục images/)
% \includegraphics[width=0.25\textwidth]{images/logo_sgu.png}\\[1.5cm]

\vspace{1.5cm}

{\LARGE\textbf{BÁO CÁO ĐỒ ÁN}}\\[1cm]

{\Large\textbf{TÊN ĐỀ TÀI:}}\\[0.5cm]

{\LARGE\textbf{GEARGO --- HỆ THỐNG CHO THUÊ ĐỒ CẮM TRẠI VÀ DU LỊCH}}\\[0.2cm]

\begin{flushleft}
\hspace{2cm}\textbf{Giảng viên hướng dẫn:}\hspace{0.5cm} Thầy Nguyễn Thanh Sang\\[0.3cm]
\hspace{2cm}\textbf{Nhóm:}\hspace{0.5cm} 10\\[0.3cm]
\hspace{2cm}\textbf{Thành viên (MSSV --- Họ tên --- Lớp):}\\[0.3cm]

\hspace{2.5cm}1. 3124560096 --- Trần Tô Thanh Tùng --- DKP1241\\[0.2cm]
\hspace{2.5cm}2. 3124560056 --- Lương Kiện Minh --- DKP1242\\[0.2cm]
\hspace{2.5cm}3. 3124560094 --- Nguyễn Minh Tú --- DKP1241\\[0.2cm]
\hspace{2.5cm}4. 3124410417 --- Trần Kim Xuyến --- DCT1246\\[0.2cm]
\hspace{2.5cm}5. 3122560038 --- Mai Trần Tuấn Kiệt --- DKP1221\\[0.2cm]
\end{flushleft}

\vfill

{\large\textbf{TP. Hồ Chí Minh, tháng 12 năm 2026}}

\end{center}
\end{titlepage}

\newpage
% =====================================================================
%  Trang NHẬN XÉT CỦA GIẢNG VIÊN
% =====================================================================
\thispagestyle{empty}

\begin{center}
{\Large\textbf{NHẬN XÉT CỦA GIẢNG VIÊN HƯỚNG DẪN}}\\[1cm]
\end{center}

\vspace{0.5cm}

\noindent
Giảng viên hướng dẫn: \textbf{Thầy Nguyễn Thanh Sang}\\[0.3cm]
Đơn vị công tác: Khoa Công nghệ Thông tin --- Trường Đại học Sài Gòn\\[0.3cm]
Nhóm sinh viên: \textbf{Nhóm 10}\\[0.3cm]
Đề tài: \textbf{GearGo --- Hệ thống cho thuê đồ cắm trại và du lịch}

\vspace{1.5cm}

\noindent
\textbf{Nội dung nhận xét:}\\[0.5cm]

\begin{spacing}{2.0}
\noindent\makebox[\linewidth]{\dotfill}

\noindent\makebox[\linewidth]{\dotfill}

\noindent\makebox[\linewidth]{\dotfill}

\noindent\makebox[\linewidth]{\dotfill}

\noindent\makebox[\linewidth]{\dotfill}

\noindent\makebox[\linewidth]{\dotfill}

\noindent\makebox[\linewidth]{\dotfill}

\noindent\makebox[\linewidth]{\dotfill}

\noindent\makebox[\linewidth]{\dotfill}

\noindent\makebox[\linewidth]{\dotfill}
\end{spacing}

\vspace{1.5cm}

\noindent
\textbf{Điểm số:}\hspace{1cm}……… / 10\hspace{1cm}\textbf{Điểm chữ:}\hspace{0.5cm}…………………

\vspace{2cm}

\begin{flushright}
TP. Hồ Chí Minh, ngày ……… tháng ……… năm 2026\\[0.3cm]
\textbf{Giảng viên hướng dẫn}\\[2cm]
Nguyễn Thanh Sang
\end{flushright}

\newpage
% =====================================================================
%  Trang LỜI CẢM ƠN
% =====================================================================
\thispagestyle{empty}

\begin{center}
{\Large\textbf{LỜI CẢM ƠN}}\\[1cm]
\end{center}

\vspace{0.5cm}

Nhóm chúng em xin gửi lời cảm ơn chân thành đến \textbf{Khoa Công nghệ Thông tin --- Trường Đại học Sài Gòn} đã tạo điều kiện cho chúng em được học tập, nghiên cứu và hoàn thành đồ án môn \textit{Công nghệ .NET} trong học kỳ này.

Đặc biệt, nhóm chúng em xin gửi lời cảm ơn sâu sắc đến thầy \textbf{Nguyễn Thanh Sang} --- giảng viên hướng dẫn, đã tận tình chỉ bảo, góp ý và đồng hành cùng nhóm trong suốt quá trình xây dựng đề tài \textbf{``GearGo --- Hệ thống cho thuê đồ cắm trại và du lịch''}. Những nhận xét của Thầy không chỉ giúp chúng em hoàn thiện sản phẩm mà còn giúp chúng em trau dồi tư duy phân tích, thiết kế hệ thống và kỹ năng làm việc nhóm.

Dù đã cố gắng hết sức, báo cáo và sản phẩm không tránh khỏi những thiếu sót. Nhóm rất mong nhận được những ý kiến đóng góp quý báu từ Thầy để sản phẩm hoàn thiện hơn và nhóm có thể tiếp tục phát triển trong các dự án tương lai.

Nhóm chúng em xin chân thành cảm ơn!

\vspace{2cm}

\begin{flushright}
TP. Hồ Chí Minh, tháng 12 năm 2026\\[0.3cm]
\textbf{Nhóm 10 --- Lập trình .NET}
\end{flushright}

\newpage

\tableofcontents
\newpage
\listoffigures
\newpage
\listoftables
\newpage

% ---------- NỘI DUNG CHƯƠNG ----------
\pagenumbering{arabic}
\setcounter{page}{1}

% =====================================================================
%  CHƯƠNG 1 — TỔNG QUAN ĐỀ TÀI & CÔNG NGHỆ
%  Người viết: Phạm Thanh Tùng (Phần 1)
%  Nguồn nội dung: Report/Assign/Phan_1_ThanhTung_TongQuan_CaiDat.md
% =====================================================================
\chapter{Tổng quan đề tài và Công nghệ}
\label{chuong:tong_quan}

% TODO: Chép nội dung từ Phan_1_ThanhTung_TongQuan_CaiDat.md — mục A

\section{Bối cảnh và lý do chọn đề tài}
\label{sec:boi_canh}

% [Nội dung 1.1 — khoảng 1-2 trang]

\section{Mục tiêu}
\label{sec:muc_tieu}

\subsection{Mục tiêu cho khách hàng (end-user)}

\subsection{Mục tiêu cho nhân viên (staff)}

\subsection{Mục tiêu cho quản trị viên (admin)}

\section{Phạm vi thực hiện}
\label{sec:pham_vi}

\subsection{Phạm vi chức năng thực hiện}

\subsection{Phạm vi không thực hiện}

\section{Công nghệ sử dụng}
\label{sec:cong_nghe}

\begin{table}[H]
\centering
\caption{Bảng công nghệ sử dụng}
\label{tab:cong_nghe}
\begin{tabularx}{\textwidth}{|l|X|}
\hline
\textbf{Thành phần} & \textbf{Công nghệ chính} \\ \hline
Backend  & ASP.NET Core Web API (.NET 10), EF Core 9, SQL Server 2022 (Docker), JWT Bearer + ASP.NET Core Identity \\ \hline
Frontend & React 19, TypeScript 5.7, Vite 8, Tailwind CSS 4, React Router DOM 7 \\ \hline
Background Job & \texttt{DonThueExpirationJob} --- IHostedService, quét hết hạn giữ chỗ mỗi 1 phút \\ \hline
Kiểm thử & xUnit (backend), Vitest (frontend), Playwright E2E (10 phân hệ frontend + backend) \\ \hline
Tài liệu & Python \texttt{python-docx} (\texttt{scripts/chuong7/}), Playwright auto-annotation (\texttt{capture\_71\_ucs.mjs}) \\ \hline
\end{tabularx}
\end{table}

\subsection{Backend --- ASP.NET Core 10}

\subsection{Frontend --- React 19 + TypeScript + Vite}

\subsection{Cơ sở dữ liệu --- SQL Server 2022 trên Docker}

\subsection{Kiểm thử tự động --- Playwright 10 phân hệ}

\section{Kế hoạch thực hiện và phân công}
\label{sec:phan_cong}

\begin{table}[H]
\centering
\caption{Bảng phân công thành viên}
\label{tab:phan_cong_thanh_vien}
\begin{tabularx}{\textwidth}{|c|l|X|c|}
\hline
\textbf{STT} & \textbf{Thành viên} & \textbf{Nhiệm vụ chính} & \textbf{\% đóng góp} \\ \hline
1 & Phạm Thanh Tùng & Lead, Auth/Infra, điều phối migration, deploy, phần báo cáo tổng quan và cài đặt & 20\% \\ \hline
2 & Kiện Minh & Catalog, tìm kiếm, khả dụng, đánh giá, thông báo, FE khách & 20\% \\ \hline
3 & Minh Tú & Khuyến mãi, kho, thiết bị, kiểm kê, logic giữ lượt & 20\% \\ \hline
4 & Kim Xuyến & Giỏ, báo giá, admin CRUD, tư vấn AI, FE admin & 20\% \\ \hline
5 & Tuấn Kiệt & Đơn thuê, thanh toán, hết hạn, bảo trì, Playwright E2E & 20\% \\ \hline
\end{tabularx}
\end{table}

% [Gantt chart tuần 1-7 — ảnh]

% =====================================================================
%  CHƯƠNG 2 — PHÂN TÍCH HỆ THỐNG
%  Người viết: Kiện Minh (Phần 2 - mục A)
%  Nguồn nội dung: Report/Assign/Phan_2_KienMinh_PhanTich_UIKhach.md
% =====================================================================
\chapter{Phân tích hệ thống}
\label{chuong:phan_tich}

\section{Đối tượng sử dụng}
\label{sec:actors}

Hệ thống GearGo phục vụ bốn đối tượng người dùng với các mức độ truy cập khác nhau.

\subsection{Khách vãng lai (Guest)}

\subsection{Khách hàng đã đăng nhập (Customer)}

\subsection{Nhân viên (Staff --- phân quyền RBAC)}

\subsection{Quản trị viên (Admin)}

\section{Danh sách chức năng theo 10 phân hệ}
\label{sec:danh_sach_chuc_nang}

\begin{table}[H]
\centering
\caption{Ma trận chức năng 10 phân hệ $\times$ 4 actors}
\label{tab:matran_chucnang}
\begin{tabularx}{\textwidth}{|l|X|c|c|c|c|}
\hline
\textbf{Phân hệ} & \textbf{Nội dung chính} & \textbf{Admin} & \textbf{Staff} & \textbf{Customer} & \textbf{Guest} \\ \hline
PH1  & Tài khoản \& Nhân sự                & X & X & X & Đăng ký \\ \hline
PH2  & Tư vấn \& Khả dụng                  & -- & X & X & X \\ \hline
PH3  & Giỏ, Giữ chỗ, Hủy                   & -- & -- & X & -- \\ \hline
PH4  & Bàn giao \& Nhận trả                & X & X & Theo dõi & -- \\ \hline
PH5  & Thanh toán, Phụ phí, Đối soát       & X & X & Thanh toán & -- \\ \hline
PH6  & Sản phẩm, Thiết bị, Bảo trì         & X & X & -- & -- \\ \hline
PH7  & NCC, Nhập kho, Kiểm kê              & X & X & -- & -- \\ \hline
PH8  & Khuyến mãi, Chính sách              & X & -- & Áp mã & -- \\ \hline
PH9  & Đánh giá \& Thông báo               & Duyệt & -- & X & -- \\ \hline
PH10 & Báo cáo \& Nhật ký                  & X & Theo quyền & -- & -- \\ \hline
\end{tabularx}
\end{table}

\section{Đặc tả chi tiết Use Case}
\label{sec:dac_ta_uc}

Phần này đặc tả 12 Use Case cốt lõi theo mẫu khung: Tên chức năng, Mục tiêu, Người thực hiện, Dữ liệu vào, Dữ liệu ra, Điều kiện trước, Luồng chính, Luồng lỗi, Quy tắc ràng buộc.

\subsection{UC1.1 --- Đăng ký tài khoản khách hàng}

\subsection{UC1.2 --- Đăng nhập hai cổng với khóa tạm 15 phút}

\subsection{UC2.1 --- Tìm kiếm và lọc sản phẩm có phân trang}

\subsection{UC2.3 --- Kiểm tra khả dụng theo lịch (Sweep Line)}

\subsection{UC3.1 --- Thêm vào giỏ thuê và tính báo giá}

\subsection{UC3.2 --- Tạo đơn thuê và giữ chỗ 15 phút}

\subsection{UC3.4 --- Hủy đơn thuê chưa thanh toán}

\subsection{UC4.2 --- Phân công mã thiết bị thực tế}

\subsection{UC4.3 --- Bàn giao thiết bị cho khách hàng}

\subsection{UC5.1 --- Thanh toán VietQR và xử lý callback idempotent}

\subsection{UC8.1 --- Áp dụng mã khuyến mãi (quy tắc giữ lượt)}

\subsection{UC9.1 --- Đánh giá sản phẩm sau khi thuê}

% =====================================================================
%  CHƯƠNG 3 — THIẾT KẾ CƠ SỞ DỮ LIỆU
%  Người viết: Minh Tú (Phần 3)
%  Nguồn nội dung: Report/Assign/Phan_3_MinhTu_CoSoDuLieu.md
% =====================================================================
\chapter{Thiết kế Cơ sở dữ liệu}
\label{chuong:co_so_du_lieu}

\section{Mô hình ERD tổng thể}
\label{sec:erd}

% \begin{figure}[H]
%   \centering
%   \includegraphics[width=\textwidth]{images/erd_tong_the.png}
%   \caption{Sơ đồ ERD toàn hệ thống GearGo --- 39 entities}
%   \label{fig:erd}
% \end{figure}

Toàn hệ thống gồm \textbf{39 entities} được chia thành 6 cụm chính:

\begin{enumerate}[leftmargin=2em]
  \item \textbf{Tài khoản và Nhân sự} (3 bảng): \texttt{TaiKhoan}, \texttt{KhachHang}, \texttt{NhanVien} + bảng quyền RBAC.
  \item \textbf{Danh mục và Sản phẩm} (3 bảng): \texttt{DanhMucSanPham}, \texttt{SanPham}, \texttt{HinhAnhSanPham}.
  \item \textbf{Kho và Thiết bị} (5 bảng): \texttt{NhaCungCap}, \texttt{PhieuNhapHang}, \texttt{ChiTietPhieuNhap}, \texttt{ThietBi}, \texttt{DieuChinhKho}.
  \item \textbf{Giỏ, Đơn và Giữ chỗ} (6 bảng): \texttt{GioThue}, \texttt{ChiTietGioThue}, \texttt{DonThue}, \texttt{ChiTietDonThue}, \texttt{GiuCho}, \texttt{LichSuTrangThaiDon}.
  \item \textbf{Khuyến mãi và Chính sách} (5 bảng): \texttt{KhuyenMai}, \texttt{KhuyenMaiSanPham}, \texttt{KhuyenMaiDanhMuc}, \texttt{LuotSuDungKhuyenMai}, \texttt{ChinhSach}.
  \item \textbf{Thanh toán, Phụ phí, Bàn giao, Bảo trì, Đánh giá, Thông báo, Nhật ký} (17 bảng còn lại).
\end{enumerate}

\section{Lược đồ cơ sở dữ liệu (RDM)}
\label{sec:rdm}

% [6 trang lược đồ tương ứng 6 cụm]

\section{Từ điển dữ liệu (Data Dictionary)}
\label{sec:data_dict}

\subsection{Bảng TAI\_KHOAN}

\subsection{Bảng KHACH\_HANG}

\subsection{Bảng SAN\_PHAM}

\subsection{Bảng THIET\_BI}

\subsection{Bảng DON\_THUE}
\label{subsec:don_thue}

\subsection{Bảng CHI\_TIET\_DON\_THUE}

\subsection{Bảng GIU\_CHO}

\subsection{Bảng LUOT\_SU\_DUNG\_KHUYEN\_MAI}

\subsection{Bảng THANH\_TOAN}

\subsection{Bảng KHUYEN\_MAI}

\subsection{Bảng CHINH\_SACH}

\subsection{Bảng PHIEU\_NHAP\_HANG}

\subsection{Bảng DANH\_GIA}

\subsection{Bảng LICH\_SU\_TRANG\_THAI\_DON}

\subsection{Bảng DIEU\_CHINH\_KHO}

\subsection{Tóm tắt các bảng phụ trợ}

\section{Quy tắc dữ liệu quan trọng}
\label{sec:quy_tac_du_lieu}

\subsection{Quy tắc ẩn / ngừng kinh doanh sản phẩm}

\subsection{Quy tắc FIFO nhập kho và phân công thiết bị}

\subsection{Quy tắc giữ chỗ 15 phút đồng bộ UTC}

\subsection{Chính sách phiên bản --- không áp hồi tố đơn cũ}

\subsection{Thuật toán tính khả dụng (Sweep Line)}

\subsection{Idempotency trong thanh toán}

\subsection{Ghi nhận lịch sử chuyển trạng thái đơn}

\subsection{Audit RBAC --- PermissionVersion}

\section{Cài đặt và Migration}
\label{sec:cai_dat_db}

% [Danh sách 7 migration EF Core + lệnh seed/reset]

% =====================================================================
%  CHƯƠNG 4 — THIẾT KẾ GIAO DIỆN & LUỒNG MÀN HÌNH
%  Người viết:
%    - Phần A (Khách vãng lai, Customer): Kiện Minh — Phan_2
%    - Phần B (Staff, Admin):             Kim Xuyến — Phan_4
% =====================================================================
\chapter{Thiết kế Giao diện và Luồng màn hình}
\label{chuong:giao_dien}

\section{Tổng quan sitemap toàn hệ thống}
\label{sec:sitemap_tong}

Hệ thống chia làm hai cổng truy cập độc lập:
\begin{itemize}
  \item \textbf{Cổng khách:} \texttt{/} (home), \texttt{/login}, \texttt{/dang-ky}, \texttt{/san-pham}, \texttt{/gio-thue}, \texttt{/don-thue}.
  \item \textbf{Cổng nội bộ:} \texttt{/staff/login}, \texttt{/staff/*} (nhân viên), \texttt{/admin/*} (quản trị).
\end{itemize}

% ---------------- PHẦN A: KHÁCH ----------------
\section{Nhánh Khách vãng lai và Khách hàng}
\label{sec:ui_khach}

\subsection{Sitemap nhánh public/customer}

\subsection{Danh sách màn hình nhánh khách}

\begin{table}[H]
\centering
\caption{Danh sách màn hình nhánh Khách (trích)}
\label{tab:screens_khach}
\begin{tabularx}{\textwidth}{|c|l|c|X|}
\hline
\textbf{Mã} & \textbf{Tên màn hình} & \textbf{Đối tượng} & \textbf{Mô tả ngắn} \\ \hline
SC-GU-01 & Trang chủ công khai         & Guest           & Hero banner, sản phẩm nổi bật \\ \hline
SC-GU-02 & Danh sách SP + lọc          & Guest/Customer  & Filter TuKhoa, DanhMuc, SucChua, Gia \\ \hline
SC-GU-03 & Chi tiết sản phẩm           & Guest/Customer  & Hình, phụ kiện, chính sách, đánh giá \\ \hline
SC-GU-04 & Kiểm tra khả dụng           & Guest/Customer  & DatePicker GioNhan/GioTra \\ \hline
SC-GU-05 & Tư vấn AI bộ đồ             & Guest/Customer  & Form nhu cầu chuyến đi \\ \hline
SC-CS-04 & Giỏ thuê                    & Customer        & Dòng SP, số lượng, lịch, áp mã \\ \hline
SC-CS-05 & Xác nhận đơn                & Customer        & Preview báo giá \& chính sách \\ \hline
SC-CS-06 & Thanh toán VietQR           & Customer        & QR + đếm ngược 15 phút \\ \hline
\end{tabularx}
\end{table}

\subsection{Screenshot và mô tả chi tiết}

% ---------------- PHẦN B: STAFF/ADMIN ----------------
\section{Nhánh Nhân viên và Quản trị viên}
\label{sec:ui_admin}

\subsection{Sitemap nhánh staff/admin}

\subsection{Danh sách màn hình nhánh staff/admin}

\begin{table}[H]
\centering
\caption{Danh sách màn hình nhánh Staff/Admin (trích)}
\label{tab:screens_admin}
\begin{tabularx}{\textwidth}{|c|l|c|X|}
\hline
\textbf{Mã} & \textbf{Tên màn hình} & \textbf{Đối tượng} & \textbf{Mô tả ngắn} \\ \hline
SC-AD-01 & Dashboard admin             & Admin  & KPI 10 phân hệ, biểu đồ doanh thu \\ \hline
SC-AD-04 & Phân quyền động             & Admin  & Matrix NV $\times$ quyền RequirePermission \\ \hline
SC-AD-09 & Chính sách phiên bản        & Admin  & Timeline bản đang hiệu lực \\ \hline
SC-ST-03 & Chuẩn bị đơn - phân công serial & Staff & Chọn mã thiết bị \\ \hline
SC-ST-04 & Phiếu bàn giao              & Staff  & Checklist, ảnh trước thuê \\ \hline
SC-ST-06 & Lập phụ phí                 & Staff  & Loại, bằng chứng, chuyển duyệt \\ \hline
SC-ST-07 & Bảng đối soát cọc           & Staff  & Thu/hoàn, chốt Hoàn tất \\ \hline
\end{tabularx}
\end{table}

\subsection{Screenshot nhánh staff/admin}

% =====================================================================
%  CHƯƠNG 5 — THIẾT KẾ XỬ LÝ & KIẾN TRÚC SOURCE
%  Người viết: Kim Xuyến (Phần 4 - mục B)
%  Nguồn nội dung: Report/Assign/Phan_4_KimXuyen_UIAdmin_KienTrucSource.md
% =====================================================================
\chapter{Thiết kế Xử lý và Kiến trúc Source}
\label{chuong:kien_truc}

\section{Kiến trúc tổng thể 3-tier}
\label{sec:kien_truc_tong_the}

Hệ thống GearGo được xây dựng theo kiến trúc 3-tier gồm:
\begin{enumerate}
  \item \textbf{Frontend React 19} --- Vite, Tailwind CSS, React Router 7, giao tiếp API qua HTTPS.
  \item \textbf{Backend ASP.NET Core 10} --- Controllers, Services, EF Core, Middleware, Background Jobs.
  \item \textbf{SQL Server 2022} --- chạy trong Docker container \texttt{geargo-sql}, migration bằng EF Core.
\end{enumerate}

\section{Kiến trúc Backend --- Layered Architecture}
\label{sec:backend_arch}

\subsection{Năm lớp chính}
\begin{itemize}
  \item \textbf{Controllers} (\texttt{backend/Controllers/}) --- HTTP routing, DTO validation, gọi service.
  \item \textbf{Services} (\texttt{backend/Services/}) --- logic nghiệp vụ, giao dịch, kiểm tra quyền.
  \item \textbf{Data} (\texttt{backend/Data/ApplicationDbContext.cs}) --- EF Core DbContext, Fluent API, seed.
  \item \textbf{Models} (\texttt{backend/Models/\{Entities, DTOs, Enums\}/}) --- entity, DTO, enum.
  \item \textbf{Cross-cutting} --- \texttt{Middleware/}, \texttt{Helpers/}, \texttt{BackgroundJobs/}, \texttt{Configuration/}.
\end{itemize}

\subsection{Quy ước đặt tên và ánh xạ SQL}

\subsection{Danh sách Controller theo domain}
\begin{itemize}
  \item \texttt{Auth/}, \texttt{Catalog/}, \texttt{Order/}, \texttt{Finance/}, \texttt{Inventory/}, \texttt{Marketing/}, \texttt{Communication/}, \texttt{Report/}, \texttt{Debug/}.
\end{itemize}

\section{Kiến trúc Frontend --- Modular}
\label{sec:frontend_arch}

\subsection{Cấu trúc thư mục \texttt{frontend/src/}}

\subsection{Quản lý trạng thái và routing}

\subsection{RBAC phía Frontend}
Component \texttt{<RequirePermission permission="..."/>} ẩn/hiện nút theo quyền động từ cache \texttt{PermissionsContext}.

\section{Thiết kế module theo khung nghiệp vụ}
\label{sec:modules}

\begin{table}[H]
\centering
\caption{Ánh xạ Module $\leftrightarrow$ Backend Service $\leftrightarrow$ Frontend Module}
\label{tab:module_mapping}
\begin{tabularx}{\textwidth}{|l|X|X|l|}
\hline
\textbf{Module} & \textbf{Backend Service} & \textbf{Controller} & \textbf{FE Module} \\ \hline
Auth           & \texttt{XacThucService}, \texttt{JwtService} & \texttt{AuthController} & \texttt{modules/auth} \\ \hline
Catalog        & \texttt{SanPhamService}, \texttt{KhaDungService} & \texttt{Catalog/*} & \texttt{modules/catalog} \\ \hline
Cart \& Quote  & \texttt{GioThueService}, \texttt{BaoGiaService} & \texttt{GioThueController} & \texttt{modules/catalog/cart} \\ \hline
Order          & \texttt{DonThueService}, \texttt{GiuChoService} & \texttt{DonThueController} & \texttt{modules/orders} \\ \hline
Payment        & \texttt{ThanhToanService} + gateway adapter & \texttt{ThanhToanController} & \texttt{modules/orders/payment} \\ \hline
Inventory      & \texttt{PhieuNhapService}, \texttt{ThietBiService} & \texttt{Inventory/*} & \texttt{modules/inventory} \\ \hline
Marketing      & \texttt{KhuyenMaiService}, \texttt{ChinhSachService} & \texttt{Marketing/*} & \texttt{modules/marketing} \\ \hline
Review+Notify  & \texttt{DanhGiaService}, \texttt{ThongBaoService} & \texttt{Communication/*} & \texttt{modules/notifications} \\ \hline
Report         & \texttt{BaoCaoService} & \texttt{Report/*} & \texttt{modules/reports} \\ \hline
\end{tabularx}
\end{table}

\section{Mười truy vấn LINQ / SQL tiêu biểu}
\label{sec:queries}

\subsection{Truy vấn 1 --- Phân trang sản phẩm công khai với lọc trạng thái kinh doanh}

\begin{lstlisting}[language=csharp, caption={LINQ --- TimKiemAsync}]
// SanPhamService.TimKiemAsync
var query = _db.SanPhams
    .Where(sp => sp.TrangThaiKinhDoanh == "DangKinhDoanh")
    .Where(sp => sp.DanhMuc.TrangThai == "HienThi");
\end{lstlisting}

\subsection{Truy vấn 2 --- Khả dụng thiết bị theo lịch (Sweep Line)}

\subsection{Truy vấn 3 --- Tạo đơn + snapshot trong transaction Serializable}

\subsection{Truy vấn 4 --- Nhập kho FIFO: phân công serial thiết bị}

\subsection{Truy vấn 5 --- Báo cáo doanh thu tách thuê / cọc / phụ phí / hoàn}

\subsection{Truy vấn 6 --- Idempotency callback thanh toán}

\subsection{Truy vấn 7 --- Job hết hạn 15 phút với UPDATE có điều kiện}

\subsection{Truy vấn 8 --- Kiểm tra mã khuyến mãi + loại lượt hết hạn}

\subsection{Truy vấn 9 --- Cập nhật quyền động và invalidate cache}

\subsection{Truy vấn 10 --- Lịch sử chuyển trạng thái đơn}

\section{Background Jobs}
\label{sec:background_jobs}

\subsection{\texttt{DonThueExpirationJob}}

\section{Cross-cutting: Exception handling và RBAC}
\label{sec:cross_cutting}

\subsection{ExceptionHandlingMiddleware}

\subsection{Attribute \texttt{[RequirePermission]}}

% =====================================================================
%  CHƯƠNG 6 — CÀI ĐẶT & TRIỂN KHAI
%  Người viết: Phạm Thanh Tùng (Phần 1 - mục B)
%  Nguồn nội dung: Report/Assign/Phan_1_ThanhTung_TongQuan_CaiDat.md
% =====================================================================
\chapter{Cài đặt và Triển khai}
\label{chuong:cai_dat}

\section{Môi trường phát triển}
\label{sec:moi_truong_dev}

\begin{itemize}
  \item \textbf{Hệ điều hành:} Windows 11 và Ubuntu 22.04 (nhóm hỗn hợp).
  \item \textbf{.NET SDK 10} --- kiểm tra bằng \texttt{dotnet --version} $\geq$ 10.
  \item \textbf{Node.js $\geq$ 20} + npm / pnpm.
  \item \textbf{Docker Desktop} --- chạy SQL Server 2022.
  \item \textbf{Python $\geq$ 3.10} --- tùy chọn, dùng script build Chương 7.
  \item \textbf{IDE:} Visual Studio 2022 / JetBrains Rider / VS Code + C\# Dev Kit.
\end{itemize}

\section{Môi trường triển khai (demo)}
\label{sec:moi_truong_deploy}

\begin{itemize}
  \item SQL Server 2022 trong Docker container \texttt{geargo-sql} --- host port \textbf{11433} $\to$ container port 1433.
  \item Backend: \url{http://localhost:5000} --- Swagger UI tại \url{http://localhost:5000/swagger}.
  \item Frontend (Vite dev server): \url{http://localhost:8443}.
\end{itemize}

\section{Hướng dẫn cài đặt và chạy}
\label{sec:huong_dan_chay}

\begin{lstlisting}[language=bash, caption={Bước 1 --- Khởi tạo SQL Server}]
docker compose up -d
\end{lstlisting}

\begin{lstlisting}[language=bash, caption={Bước 2 --- Chạy Backend (tự migrate + seed)}]
dotnet restore backend/GearGo.csproj
dotnet run --project backend/GearGo.csproj --urls http://localhost:5000
\end{lstlisting}

\begin{lstlisting}[language=bash, caption={Bước 3 --- Chạy Frontend}]
npm --prefix frontend install
npm --prefix frontend run dev -- --port 8443 --host
\end{lstlisting}

\section{Tài khoản seed và hai cổng đăng nhập}
\label{sec:tai_khoan_seed}

\begin{itemize}
  \item \textbf{Cổng Khách hàng:} \url{http://localhost:8443/login}
  \item \textbf{Cổng Nhân viên / Quản trị:} \url{http://localhost:8443/staff/login}
  \item \textbf{Mật khẩu seed chung:} \texttt{GearGo@2026}
\end{itemize}

\begin{table}[H]
\centering
\caption{Tài khoản seed phục vụ chấm điểm}
\label{tab:tai_khoan_seed}
\begin{tabularx}{\textwidth}{|l|l|l|X|}
\hline
\textbf{Cổng} & \textbf{Email / SĐT} & \textbf{Vai trò} & \textbf{Ghi chú} \\ \hline
Khách hàng  & customer@geargo.vn (0901000003) & KhachHang    & Nguyễn Văn An \\ \hline
Khách hàng  & quang.kh@geargo.vn               & KhachHang    & Trần Minh Quang \\ \hline
Khách hàng  & locked.kh@geargo.vn              & KhachHang    & Tài khoản bị khóa --- test UC1.2 \\ \hline
Nội bộ      & admin@geargo.vn (0901000001)    & QuanTriVien  & Nguyễn Văn Quản Trị (Full quyền) \\ \hline
Nội bộ      & staff@geargo.vn (0901000002)    & NhanVien     & Lê Thị Mai \\ \hline
Nội bộ      & tuan.staff@geargo.vn (0901000005) & NhanVien   & Trần Tuấn \\ \hline
\end{tabularx}
\end{table}

\section{Reset và seed lại cơ sở dữ liệu}
\label{sec:reset_db}

\begin{lstlisting}[language=bash]
dotnet ef database drop -f --project backend/GearGo.csproj
dotnet run --project backend/GearGo.csproj --urls http://localhost:5000
\end{lstlisting}

\section{Cấu hình quan trọng}
\label{sec:cau_hinh}

\begin{itemize}
  \item \texttt{backend/appsettings.json} --- connection string, JWT secret, cấu hình VietQR.
  \item \texttt{frontend/.env.development} --- \texttt{VITE\_API\_URL=http://localhost:5000}, \texttt{VITE\_USE\_MOCK=false}.
  \item Không hard-code đường dẫn tuyệt đối theo máy; mọi path tương đối từ root repo.
\end{itemize}

% =====================================================================
%  CHƯƠNG 7 — KIỂM THỬ
%  Người viết: Tuấn Kiệt (Phần 5 - mục A)
%  Nguồn nội dung: Report/Assign/Phan_5_TuanKiet_KiemThu_KetLuan.md
% =====================================================================
\chapter{Kiểm thử}
\label{chuong:kiem_thu}

\section{Chiến lược kiểm thử ba tầng}
\label{sec:chien_luoc_test}

\begin{enumerate}
  \item \textbf{Unit test (xUnit)} --- \texttt{backend/test/GearGo.Tests/} --- kiểm thử domain service độc lập.
  \item \textbf{Component test (Vitest)} --- \texttt{frontend/} --- kiểm thử React component và custom hook.
  \item \textbf{End-to-end test (Playwright)} --- \texttt{tests/playwright/} --- kiểm thử luồng thật qua browser và API thật.
\end{enumerate}

\section{Playwright E2E --- 10 phân hệ $\times$ 2 nhánh}
\label{sec:playwright}

Bộ test Playwright gồm:
\begin{itemize}
  \item 10 spec \textbf{frontend} UI: \texttt{tests/playwright/frontend/phan\_he\_XX\_*/phan\_he\_XX.spec.ts}
  \item 10 spec \textbf{backend} API: \texttt{tests/playwright/backend/phan\_he\_XX\_*/phan\_he\_XX.spec.ts}
\end{itemize}

\begin{lstlisting}[language=bash, caption={Lệnh chạy Playwright}]
npm --prefix tests/playwright run test:frontend
npm --prefix tests/playwright run test:backend
npm --prefix tests/playwright test
\end{lstlisting}

Mỗi test tự seed dữ liệu và cleanup sau khi hoàn tất.

\section{Bảng Test Cases (22 TC)}
\label{sec:tc}

\begin{longtable}{|c|c|p{3.5cm}|p{4cm}|p{2.5cm}|p{3cm}|c|}
\caption{Bảng Test Cases tiêu biểu} \label{tab:test_cases} \\
\hline
\textbf{TC} & \textbf{PH} & \textbf{Chức năng} & \textbf{Bước thực hiện} & \textbf{Dữ liệu} & \textbf{KQ mong đợi} & \textbf{KQ} \\ \hline
\endfirsthead
\multicolumn{7}{c}{{\bfseries \tablename\ \thetable{} --- tiếp}} \\
\hline
\textbf{TC} & \textbf{PH} & \textbf{Chức năng} & \textbf{Bước thực hiện} & \textbf{Dữ liệu} & \textbf{KQ mong đợi} & \textbf{KQ} \\ \hline
\endhead

TC-01 & PH1 & Đăng ký khách thành công        & Mở \texttt{/dang-ky}, submit & test@geargo.vn / Abc@12345 & 201, redirect \texttt{/login} & Pass \\ \hline
TC-02 & PH1 & Khóa 15' sau 5 lần sai          & Login sai 5 lần              & mật khẩu sai & 429 Tạm khóa 15p & Pass \\ \hline
TC-03 & PH1 & Cổng nội bộ từ chối role khách  & \texttt{/staff/login}        & customer@geargo.vn & 403 & Pass \\ \hline
TC-04 & PH2 & Tìm kiếm SP phân trang          & GET \texttt{/san-pham}       & page=1, size=12 & 12 SP, totalPages > 1 & Pass \\ \hline
TC-05 & PH2 & Khả dụng Sweep Line             & GET /san-pham/\{id\}         & 2 lịch giao nhau & Tổng - max chồng & Pass \\ \hline
TC-06 & PH2 & Ẩn SP \texttt{TamNgung}         & GET public list              & -- & Không xuất hiện & Pass \\ \hline
TC-07 & PH3 & Cộng dồn SL khi thêm 2 lần      & POST /gio-thue/them $\times$2 & SL 1 + 2 & 1 dòng, SL = 3 & Pass \\ \hline
TC-08 & PH3 & Mã \texttt{TEST10} giảm 10\%    & Áp mã                        & tiền thuê 600k & giảm 60k & Pass \\ \hline
TC-09 & PH3 & Tạo đơn + GiuCho hạn 15'        & POST /don-thue              & giỏ 3 dòng & 3 GiuCho DangGiu & Pass \\ \hline
TC-10 & PH3 & Hủy đơn giải phóng giữ chỗ      & POST /don-thue/\{id\}/huy   & -- & GiuCho DaGiaiPhong & Pass \\ \hline
TC-11 & PH3 & Job tự HetHan sau 15'           & Chờ > 60s                   & -- & DonThue HetHan & Pass \\ \hline
TC-12 & PH4 & Phân công serial thiết bị      & Staff Chuẩn bị đơn          & -- & ChiTietDonThue có MaThietBi & Pass \\ \hline
TC-13 & PH4 & Nhận trả một phần               & Submit 2/3 thiết bị         & -- & Đơn vẫn DangThue + nhãn Trả một phần & Pass \\ \hline
TC-14 & PH5 & Thanh toán mock $\to$ DaXacNhan & POST /tao-url $\to$ callback & -- & 2 ChiTietThanhToan (Thuê + Cọc) & Pass \\ \hline
TC-15 & PH5 & Idempotent MaYeuCau             & Gọi callback 2 lần          & cùng MaYeuCau & Không ghi trùng & Pass \\ \hline
TC-16 & PH5 & Admin duyệt phụ phí tranh chấp  & Staff lập $\to$ Admin duyệt & -- & TrangThai DaDuyet & Pass \\ \hline
TC-17 & PH6 & Thay quyền $\to$ invalidate cache & Admin thu hồi $\to$ staff gọi API & -- & 403 ngay lần kế & Pass \\ \hline
TC-18 & PH7 & Nhập kho sinh ThietBi          & Admin xác nhận PhieuNhap    & -- & N ThietBi SanSang & Pass \\ \hline
TC-19 & PH8 & Mã KM vượt tổng lượt           & Khách N+1 áp mã giới hạn N  & -- & 400 KhuyenMaiKhongHopLe & Pass \\ \hline
TC-20 & PH8 & Lượt hết hạn không chiếm lượt  & Đơn hết hạn chưa trả        & -- & Khách kế tiếp dùng được mã & Pass \\ \hline
TC-21 & PH9 & 1 SP / 1 đơn chỉ đánh giá 1 lần & POST /danh-gia $\times$2    & -- & Lần 2 $\to$ 409 Conflict & Pass \\ \hline
TC-22 & PH10 & Báo cáo doanh thu tách 4 cột   & GET /bao-cao/doanh-thu      & kỳ tùy chọn & 4 cột Thuê/Cọc/Phụ phí/Hủy & Pass \\ \hline
\end{longtable}

\section{Race condition và concurrency tests}
\label{sec:race}

\begin{itemize}
  \item \textbf{Hai user cùng đặt thiết bị cuối} $\to$ chỉ một thành công (Isolation Serializable + sweep line check).
  \item \textbf{Thanh toán chạy cùng job hết hạn} $\to$ không ghi đè (UPDATE có điều kiện \texttt{WHERE TrangThai = 'ChoThanhToan'}).
  \item \textbf{Callback gửi trùng} $\to$ idempotent nhờ \texttt{ThanhToan.MaYeuCau} UNIQUE.
\end{itemize}

\section{Backend Unit test (xUnit)}
\label{sec:unit_test}

\begin{lstlisting}[language=bash]
dotnet test GearGo_Webapp.sln
\end{lstlisting}

\section{Kết quả tổng hợp}
\label{sec:ket_qua_test}

% =====================================================================
%  KẾT LUẬN & HƯỚNG PHÁT TRIỂN
%  Người viết: Tuấn Kiệt (Phần 5 - mục B)
% =====================================================================
\chapter*{KẾT LUẬN VÀ HƯỚNG PHÁT TRIỂN}
\addcontentsline{toc}{chapter}{Kết luận và Hướng phát triển}

\section*{Kết quả đạt được}
\addcontentsline{toc}{section}{Kết quả đạt được}

Sau 7 tuần triển khai, nhóm đã hoàn thành \textbf{10 phân hệ nghiệp vụ} với \textbf{71 màn hình} và \textbf{39 entities}, đáp ứng đầy đủ các yêu cầu của đề tài:

\begin{itemize}
  \item \textbf{PH1 --- Tài khoản:} hai cổng đăng nhập độc lập, RBAC động \texttt{[RequirePermission]}, khóa tạm 15 phút sau 5 lần sai, phân quyền nhân viên linh động.
  \item \textbf{PH2 --- Tư vấn và Khả dụng:} thuật toán Sweep Line, lọc 7 tiêu chí, tư vấn AI mock.
  \item \textbf{PH3 --- Giỏ, Giữ chỗ, Hủy:} giỏ cộng dồn, giữ chỗ 15 phút đồng bộ UTC, áp mã khuyến mãi.
  \item \textbf{PH4 --- Bàn giao và Nhận trả:} phân công serial, phiếu bàn giao, nhận trả một phần.
  \item \textbf{PH5 --- Thanh toán và Đối soát:} VietQR 15 phút, callback idempotent, phụ phí và duyệt tranh chấp.
  \item \textbf{PH6--7 --- Vòng đời thiết bị và Nhập kho:} FIFO, kiểm kê, điều chỉnh kho.
  \item \textbf{PH8 --- Khuyến mãi và Chính sách:} chính sách phiên bản, không áp hồi tố.
  \item \textbf{PH9 --- Đánh giá và Thông báo:} đánh giá sau thuê, phản hồi cửa hàng, thông báo deep link.
  \item \textbf{PH10 --- Báo cáo và Nhật ký:} tách doanh thu thuê/cọc/phụ phí/hủy, audit log.
\end{itemize}

\section*{Hạn chế}
\addcontentsline{toc}{section}{Hạn chế}

\begin{itemize}
  \item Chưa gửi email thật (chỉ log console), SendGrid chưa bật production.
  \item VietQR sandbox --- chưa ký hợp đồng cổng thanh toán thật.
  \item Chatbot AI (n8n + Gemini) đã có skeleton nhưng chưa deploy cloud.
  \item Chưa thực hiện load test với 1000 user đồng thời.
  \item Mobile responsive dưới 375px chưa polish hoàn thiện.
\end{itemize}

\section*{Hướng phát triển}
\addcontentsline{toc}{section}{Hướng phát triển}

\begin{enumerate}
  \item Tích hợp cổng thanh toán production: VietQR, VNPay, MoMo.
  \item SendGrid email + Firebase Cloud Messaging cho push notification.
  \item Triển khai chatbot n8n + Gemini production, kèm RAG trên knowledge base nghiệp vụ.
  \item Phân tích ML: dự báo nhu cầu theo mùa, gợi ý khuyến mãi động theo lịch sử khách.
  \item Ứng dụng di động React Native.
  \item Multi-tenant: hỗ trợ nhiều chi nhánh, đồng bộ kho liên chi nhánh.
\end{enumerate}

% =====================================================================
%  PHỤ LỤC
%    A - Thanh Tùng (git, deploy, config)
%    B - Tuấn Kiệt (71 UC screenshots, SQL seed, Playwright report)
% =====================================================================
\appendix

\chapter{Git, Cấu hình và Hướng dẫn chấm điểm}
\label{apx:a}

\section{Link mã nguồn}
\url{https://github.com/ThanhTung0401/GearGo_Webapp}

\section{Cấu trúc thư mục dự án}

\section{Danh sách EF Core Migration}
\begin{enumerate}
  \item \texttt{20260924045627\_FullERDInitialCreate}
  \item \texttt{20260924060134\_RemoveRedundantDonThueForeignKey}
  \item \texttt{20261110084459\_AddPermissionIdentity}
  \item \texttt{20261113163000\_AddDanhGiaPhanHoiCuaHang}
  \item \texttt{20261114062923\_AddIdempotencyAndQuaHanContactLog}
  \item \texttt{20261119062659\_AddTaiKhoanPermissionVersion}
\end{enumerate}

\section{Docker compose}

\section{Phân công chi tiết theo tuần}

\chapter{Bộ 71 Màn hình Use Case và Dữ liệu seed}
\label{apx:b}

\section{Danh mục 71 UC screenshots}

\section{Script SQL seed đầy đủ}

\section{Playwright HTML report}

\section{Video demo}


\end{document}
